Existing hardware-aware NAS (HW-NAS) methods typically assume access to precise information about the target device, either via analytical approximations of post-compilation latency or through learned latency predictors. Such approaches limit generalization and can introduce estimation errors that are detrimental in risk-sensitive applications.

We propose a two-stage HW-NAS framework, in which we first \emph{meta-learn} an architecture controller on a distribution of synthetic devices. At test time, the trained controller deploys directly to the target device without relying on any pre-collected information, and rather only exploits direct interactions. In particular, the controller designs an architecture for the target device by interacting with it through a small number of high-fidelity latency measurements, thanks to the meta-learning phase. Thus, our approach dispenses with amortized latency estimation at deployment, reducing model-error risks and enabling robust HW-NAS in scenarios where one cannot trust approximate latency sources.

To control computational costs, we train our controller with training-free accuracy proxies, allowing us to scale the meta-training phase without incurring the overhead of full network training. We benchmark on HW-NATS-Bench, demonstrating that our method \emph{generalizes to unseen devices} and searches for latency-efficient architectures by in-context adaptation using only a few real-world latency evaluations at test-time.